\documentclass[12pt,a4paper]{report}

\usepackage[T1]{fontenc}
\usepackage[utf8]{inputenc}
\usepackage[french]{babel}
\usepackage{geometry}
\geometry{margin=2.5cm}

\usepackage{graphicx}
\usepackage{booktabs}
\usepackage{longtable}
\usepackage{tabularx}
\usepackage{float}
\usepackage{hyperref}
\hypersetup{
  colorlinks=true,
  linkcolor=blue,
  urlcolor=blue,
  citecolor=blue
}

\usepackage{amsmath}
\usepackage{amssymb}
\usepackage{mathtools}

\usepackage{xcolor}
\usepackage{listings}
\lstset{
  basicstyle=\small\ttfamily,
  breaklines=true,
  frame=single,
  numbers=left,
  numberstyle=\tiny,
  showstringspaces=false
}

\usepackage{tikz}
\usetikzlibrary{positioning,arrows.meta,shapes.geometric}

\usepackage{enumitem}
\setlist[itemize]{label={},leftmargin=*,itemsep=0.2em,topsep=0.2em}
\setlist[enumerate]{leftmargin=*,itemsep=0.2em,topsep=0.2em}

\title{Cahier de conception\\\large Plateforme Comptabilité --- ERP Yowyob}
\author{Auteur : AZANGUE LEONEL DELMAT\\Superviseur : Pr Dr Eng Djotio Thomas NDIE}
\date{\today}

\begin{document}
\maketitle
\tableofcontents
\listoffigures

\chapter{Introduction}
\section{Objectif}
Ce document est un \textbf{cahier de conception complet} de la \textbf{plateforme comptabilité} de l'ERP Yowyob.
Il vise à décrire l'architecture, les modèles de données, les flux applicatifs, l'API, la sécurité et le multi-tenant, ainsi qu'une \textbf{modélisation mathématique} des invariants comptables (double entrée, contraintes de périodes, lettrage, pointage).

\section{Périmètre}
Le périmètre porte principalement sur le module \texttt{accounting} (comptabilité générale OHADA) et les composants transverses \texttt{common} et \texttt{config}.

\chapter{Vue d'ensemble du système}
\section{Stack technique}
\begin{table}[H]
\centering
\begin{tabularx}{\textwidth}{@{}lX@{}}
\toprule
\textbf{Composant} & \textbf{Choix technologique} \\
\midrule
Langage / Runtime & Java 21 \\
Framework HTTP & Spring Boot 3.5.3, Spring WebFlux (réactif) \\
Base de données & PostgreSQL, accès réactif via R2DBC \\
Gestion de schéma & Liquibase (exécuté au démarrage) \\
Sécurité & Spring Security WebFlux + JWT (validation via service externe) \\
Composants optionnels & Redis (cache), Kafka (événements/audit), Elasticsearch (recherche) \\
\bottomrule
\end{tabularx}
\caption{Stack technique de la plateforme comptabilité}
\end{table}

\section{Organisation du code}
Le code est organisé en trois ensembles cohérents:
\begin{enumerate}
  \item \textbf{\texttt{com.yowyob.erp.accounting}}: domaine comptable (API, services, persistance, entités, DTO).
  \item \textbf{\texttt{com.yowyob.erp.common}}: éléments transverses (wrappers de réponse, exceptions, objets communs).
  \item \textbf{\texttt{com.yowyob.erp.config}}: configuration technique (auth/JWT, CORS, multi-organisation, R2DBC, Redis, Kafka, Swagger).
\end{enumerate}

\section{Architecture logique (couches)}
\begin{figure}[H]
\centering
\begin{tikzpicture}[
  layer/.style={draw, rounded corners, minimum width=13cm, minimum height=1cm, align=center},
  arrow/.style={-Latex, thick}
]
\node[layer, fill=blue!10] (api) {Couche API \\ \texttt{controller} (REST WebFlux)};
\node[layer, fill=blue!10, below=0.6cm of api] (svc) {Couche Métier \\ \texttt{service} (règles, validations, workflows)};
\node[layer, fill=blue!10, below=0.6cm of svc] (repo) {Couche Persistance \\ \texttt{repository} (R2DBC, requêtes SQL)};
\node[layer, fill=blue!10, below=0.6cm of repo] (db) {Stockage \\ PostgreSQL (schéma versionné par Liquibase)};
\draw[arrow] (api) -- (svc);
\draw[arrow] (svc) -- (repo);
\draw[arrow] (repo) -- (db);
\end{tikzpicture}
\caption{Architecture en couches (vue reverse-engineered)}
\end{figure}

\chapter{Exigences (extraites du code et du README)}
\section{Exigences fonctionnelles}
\begin{enumerate}[label=\textbf{RF-\arabic*.}]
  \item Gestion du plan comptable OHADA, des comptes, journaux, écritures et détails d'écritures.
  \item Gestion des exercices et périodes comptables, avec mécanismes de clôture.
  \item Lettrage et pointage des lignes d'écritures.
  \item Journalisation (audit) des actions et traçabilité des changements.
  \item Génération d'écritures à partir d'objets métier comptabilisables (\texttt{ComptableObject}).
\end{enumerate}

\section{Exigences non fonctionnelles}
\begin{enumerate}[label=\textbf{RNF-\arabic*.}]
  \item Réactivité et scalabilité (WebFlux + R2DBC) pour supporter des volumes importants d'écritures.
  \item Isolation multi-organisation: toutes les données métier sont partitionnées par \texttt{organization\_id}.
  \item Traçabilité: journal d'audit et champs d'audit (\texttt{created\_at}, \texttt{updated\_at}, \texttt{created\_by}, \texttt{updated\_by}).
  \item Sécurité: authentification JWT, contrôle d'accès par rôles, CORS configurable.
\end{enumerate}

\chapter{Multi-tenant et isolation des données}
\section{Principe}
L'isolation des données est réalisée par propagation d'un identifiant d'organisation (\texttt{organizationId}) dans le \textbf{Reactor Context}.
Chaque requête API est associée à une organisation déterminée via un header configurable (par défaut \texttt{X-Tenant-ID}) ou un paramètre de requête \texttt{tenantId} (cas SSE).

\section{Résolution du tenant}
Un \texttt{WebFilter} (\texttt{OrganizationWebFilter}) extrait l'identifiant et l'injecte dans le contexte réactif.
\subsection{Règles de résolution}
\begin{enumerate}
  \item Lire le header \texttt{X-Tenant-ID} (ou nom configuré).
  \item Sinon, lire le query param \texttt{tenantId}.
  \item Sinon, utiliser une valeur par défaut (\texttt{app.organization.default-organization}).
\end{enumerate}

\chapter{Sécurité}
\section{Authentification}
L'authentification repose sur un token \texttt{Bearer} dans \texttt{Authorization}. Un filtre (\texttt{JwtAuthenticationFilter}) appelle un service \texttt{AuthService} pour valider le token via une API externe (\texttt{auth.api.url}).
En cas de succès, il place une authentification avec des autorités \texttt{ROLE\_X} dans le contexte de sécurité réactif.

\section{Autorisation}
La chaîne de sécurité (\texttt{SecurityWebFilterChain}) définit des règles \texttt{pathMatchers}.
\textbf{Point critique}: une règle de test autorise \texttt{/api/accounting/**} sans token (mentionnée comme TEMP dans le code).

\section{Gestion des erreurs}
Les exceptions sont normalisées via \texttt{GlobalExceptionHandler} et retournées sous une forme \texttt{ApiResponseWrapper}.

\chapter{Conception des données}
\section{Stratégie de migration}
Le schéma SQL est versionné par Liquibase et orchestré depuis \texttt{db.changelog-master.xml}.
Une migration importante (\texttt{changeset-017-unify-tenant-org.xml}) unifie les notions \texttt{tenant} et \texttt{organization} et renomme les colonnes \texttt{tenant\_id} en \texttt{organization\_id}.

\section{Entités principales (comptabilité)}
Ce projet utilise des entités R2DBC annotées \texttt{@Table}.
Exemples représentatifs:
\begin{itemize}
  \item \texttt{EcritureComptable} \(\rightarrow\) \texttt{ecritures\_comptables}
  \item \texttt{DetailEcriture} \(\rightarrow\) \texttt{details\_ecritures}
  \item \texttt{Compte} \(\rightarrow\) \texttt{comptes}
  \item \texttt{JournalComptable} \(\rightarrow\) \texttt{journaux\_comptables}
  \item \texttt{PeriodeComptable} \(\rightarrow\) \texttt{periodes\_comptables}
  \item \texttt{JournalAudit} \(\rightarrow\) \texttt{journal\_audit}
\end{itemize}

\section{Modèle relationnel (extrait)}
\begin{longtable}{@{}p{4cm}p{10cm}@{}}
\toprule
\textbf{Table} & \textbf{Rôle} \\
\midrule
\endhead
\texttt{organizations} & Organisation (tenant). Identité et métadonnées (nom, email, téléphone, tax\_id, etc.). \\
\texttt{agences} & Agences rattachées à une organisation. \\
\texttt{comptes} & Comptes comptables (classe OHADA, type, solde). \\
\texttt{journaux\_comptables} & Journaux (achats, ventes, trésorerie, etc.). \\
\texttt{exercices\_comptables} & Exercice fiscal/comptable (dates, clôture). \\
\texttt{periodes\_comptables} & Périodes (liées à un exercice), ouverture/clôture. \\
\texttt{ecritures\_comptables} & Entête d'écriture comptable (numéro, libellé, totaux, statut, validée). \\
\texttt{details\_ecritures} & Lignes de débit/crédit par compte (lettrage, pointage). \\
\texttt{journal\_audit} & Journalisation des actions (création, validation, suppression...). \\
\bottomrule
\end{longtable}

\chapter{API REST (reverse engineering) --- Comptabilité}
\section{Conventions}
Les endpoints sont regroupés sous le préfixe \texttt{/api/accounting/} (ex: \texttt{/ecritures}, \texttt{/comptes}, \texttt{/journaux}, etc.).
Les contrôleurs sont annotés \texttt{@RestController} et retournent des \texttt{Mono<ResponseEntity<ApiResponseWrapper<...>>>}.

\section{Exemple: \texttt{/api/accounting/ecritures}}
Le contrôleur \texttt{EcritureComptableController} expose notamment:
\begin{itemize}
  \item \texttt{POST /api/accounting/ecritures} --- création
  \item \texttt{PUT /api/accounting/ecritures/\{id\}} --- mise à jour
  \item \texttt{POST /api/accounting/ecritures/\{id\}/validate} --- validation
  \item \texttt{GET /api/accounting/ecritures} --- liste
  \item \texttt{GET /api/accounting/ecritures/non-validated} --- non validées
  \item \texttt{GET /api/accounting/ecritures/search} --- recherche
  \item \texttt{POST /api/accounting/ecritures/generate} --- génération depuis \texttt{ComptableObject}
  \item \texttt{PATCH /api/accounting/ecritures/\{id\}/cancel} --- annulation
  \item \texttt{PATCH /api/accounting/ecritures/\{id\}/deactivate} --- désactivation
  \item \texttt{DELETE /api/accounting/ecritures/\{id\}} --- suppression
\end{itemize}

\section{Couverture des contrôleurs du module comptable}
Le module comptabilité expose (au moment du reverse engineering) les contrôleurs suivants:
\begin{itemize}
  \item \texttt{CompteController} \(\rightarrow\) \texttt{/api/accounting/comptes}
  \item \texttt{PlanComptableController} \(\rightarrow\) \texttt{/api/accounting/plan-comptable}
  \item \texttt{JournalComptableController} \(\rightarrow\) \texttt{/api/accounting/journals}
  \item \texttt{EcritureComptableController} \(\rightarrow\) \texttt{/api/accounting/ecritures}
  \item \texttt{PeriodeComptableController} \(\rightarrow\) \texttt{/api/accounting/periodes}
  \item \texttt{ExerciceComptableController} \(\rightarrow\) \texttt{/api/accounting/exercices}
  \item \texttt{RapportController} \(\rightarrow\) \texttt{/api/accounting/rapport}
  \item \texttt{PointageController} \(\rightarrow\) \texttt{/api/accounting/pointage}
  \item \texttt{LettrageController} \(\rightarrow\) \texttt{/api/comptable/lettrage} (préfixe atypique)
  \item Autres contrôleurs du domaine (\texttt{taxes}, \texttt{devises}, \texttt{immobilisations}, \texttt{budgets}, \texttt{paie}, \texttt{brouillards}, etc.)
\end{itemize}

\section{Annexe API exhaustive}
La liste exhaustive des endpoints (tous contrôleurs, méthodes, chemins, verbes HTTP) est fournie en annexe.

\chapter{Modélisation mathématique (à ne pas négliger)}
\section{Notations}
Soit une organisation \(o \in \mathcal{O}\). Une écriture comptable \(e\) est un couple:
\[
e = (H, L)
\]
où \(H\) est l'entête (journal, période, date, statut) et \(L\) est l'ensemble des lignes \(l \in L\).

Chaque ligne \(l\) est associée à un compte \(a \in \mathcal{A}\) et porte des montants \((d_l, c_l)\) (débit, crédit) avec \(d_l \ge 0\), \(c_l \ge 0\).

\section{Invariant de la double entrée}
Une écriture est \textbf{équilibrée} si:
\[
\sum_{l \in L} d_l = \sum_{l \in L} c_l
\]
Dans l'implémentation, un seuil de tolérance \(\varepsilon = 0{,}01\) est utilisé:
\[
\left|\sum d_l - \sum c_l\right| \le \varepsilon
\]
Cette contrainte est vérifiée lors de la création/mise à jour/validation (\texttt{validateBalance}).

\section{Balance, grand livre, compte de résultat, bilan (formalisme)}
\subsection{Grand livre (par compte)}
Soit un compte \(a\) et une période \([t_0,t_1]\). On définit l'ensemble des lignes:
\[
L(a,t_0,t_1)=\{\,l \mid \mathrm{account}(l)=a \wedge t_0 \le \mathrm{date}(l) \le t_1\,\}
\]
Les totaux de mouvements sont:
\[
D(a)=\sum_{l \in L(a,t_0,t_1)} d_l,\quad C(a)=\sum_{l \in L(a,t_0,t_1)} c_l
\]
Le solde de clôture s'exprime à partir d'un solde d'ouverture \(S_0(a)\):
\[
S_1(a)=S_0(a)+D(a)-C(a)
\]
(la convention de signe peut varier selon la nature du compte; le système peut aussi porter des soldes débit/crédit séparés).

\subsection{Balance des comptes}
La balance est une agrégation sur un ensemble de comptes \(\mathcal{A}\) produisant, pour chaque \(a\), un quadruplet:
\[
\big(S_0(a), D(a), C(a), S_1(a)\big)
\]
La contrainte de cohérence globale (si tout est saisi en double entrée) est:
\[
\sum_{a \in \mathcal{A}} D(a) = \sum_{a \in \mathcal{A}} C(a)
\]

\subsection{Compte de résultat}
Soit \(\mathcal{A}_{\text{prod}}\) l'ensemble des comptes de produits (classe 7, selon paramétrage), et \(\mathcal{A}_{\text{ch}}\) l'ensemble des comptes de charges (classe 6).
On définit:
\[
\mathrm{PRODUITS}=\sum_{a \in \mathcal{A}_{\text{prod}}} (C(a)-D(a)),\quad
\mathrm{CHARGES}=\sum_{a \in \mathcal{A}_{\text{ch}}} (D(a)-C(a))
\]
Le résultat net sur la période est:
\[
\mathrm{RN}=\mathrm{PRODUITS}-\mathrm{CHARGES}
\]

\subsection{Bilan}
Soit \(\mathcal{A}_{\text{actif}}\) et \(\mathcal{A}_{\text{passif}}\) (classes 1--5 selon mapping OHADA et paramétrage).
On agrège les soldes:
\[
\mathrm{ACTIF}=\sum_{a \in \mathcal{A}_{\text{actif}}} \max(S_1(a),0),\quad
\mathrm{PASSIF}=\sum_{a \in \mathcal{A}_{\text{passif}}} \max(-S_1(a),0)
\]
et une propriété attendue est l'équilibre bilan:
\[
\mathrm{ACTIF} = \mathrm{PASSIF}
\]
La mise en œuvre concrète dépend des conventions de signe et de la classification des comptes.

\section{Contraintes de période}
Soit une période \(p\) avec un indicateur \(\mathrm{closed}(p) \in \{0,1\}\).
La règle métier est:
\[
\mathrm{closed}(p)=1 \Rightarrow \neg \mathrm{create}(e) \wedge \neg \mathrm{update}(e)
\]
c'est-à-dire: aucune écriture ne peut être créée/modifiée dans une période clôturée.

\section{Machine à états (écritures)}
Soit \(S = \{\texttt{BROUILLON}, \texttt{VALIDE}, \texttt{ANNULE}\}\).
Transitions autorisées (extrait des services):
\[
\texttt{BROUILLON} \xrightarrow{\mathrm{validate}} \texttt{VALIDE}, \quad
\texttt{BROUILLON} \xrightarrow{\mathrm{cancel}} \texttt{ANNULE}, \quad
\texttt{VALIDE} \xrightarrow{\mathrm{cancel}} \texttt{ANNULE}
\]
Avec contrainte: \(\mathrm{validate}\) n'est possible que si l'écriture est équilibrée.

\section{Lettrage}
Chaque ligne peut être lettrée (\texttt{lettree}).
Une formalisation possible consiste à définir une relation de lettrage \(\mathcal{R}\) sur les lignes d'un même compte:
\[
\mathcal{R} \subseteq L \times L
\]
et imposer qu'un groupe lettré \(G \subseteq L\) respecte l'équilibre partiel:
\[
\sum_{l \in G} d_l = \sum_{l \in G} c_l
\]
Le code stocke actuellement l'état de lettrage au niveau des lignes (\texttt{lettree}, \texttt{date\_lettrage}); les règles de groupement peuvent être portées par un service de lettrage.

\section{Pointage bancaire}
Pour une ligne \(l\), le pointage est \(pointee(l) \in \{0,1\}\) avec une référence bancaire \(\mathrm{ref}(l)\).
Une contrainte d'unicité métier recommandée est:
\[
\forall l_1 \ne l_2,\ \mathrm{ref}(l_1)=\mathrm{ref}(l_2) \wedge pointee(l_1)=pointee(l_2)=1 \Rightarrow \text{invalide}
\]
(à implémenter éventuellement via contrainte BD ou règle applicative selon le besoin).

\chapter{Conception des services (extrait)}
\section{Service des écritures}
Le service \texttt{EcritureComptableService} illustre les règles métier:
\begin{itemize}
  \item création d'une écriture \(\rightarrow\) génération d'un numéro \texttt{ECR-xxxxxxxx}
  \item contrôle période non clôturée
  \item validation d'équilibre débit/crédit
  \item journalisation (table audit) + émission d'événement
  \item cache Redis (liste, recherche) avec invalidation lors des mutations
\end{itemize}

\chapter{Déploiement et exploitation}
\section{Configuration}
La configuration est centralisée dans \texttt{application.properties} et surchargée par variables d'environnement:
\begin{itemize}
  \item DB: \texttt{DB\_HOST}, \texttt{DB\_PORT}, \texttt{DB\_NAME}, \texttt{POSTGRES\_USER}, \texttt{POSTGRES\_PASSWORD}
  \item Auth: \texttt{AUTH\_API\_URL}, \texttt{JWT\_SECRET}
  \item Features: \texttt{SPRING\_KAFKA\_ENABLED}, \texttt{SPRING\_ELASTICSEARCH\_ENABLED}, \texttt{SPRING\_CACHE\_TYPE}
\end{itemize}

\section{Docker / Compose}
\begin{itemize}
  \item \texttt{docker-compose.yml} démarre Postgres, Redis, Kafka, Elasticsearch et l'application.
  \item \texttt{Dockerfile} compile et exécute l'application.
\end{itemize}

\section{CI/CD}
Le dépôt contient un workflow GitHub Actions pour construire et publier une image Docker.

\chapter{Risques, dette technique et recommandations}
\begin{itemize}
  \item \textbf{Sécurité}: routes comptables temporairement \texttt{permitAll} \(\rightarrow\) à corriger avant production.
  \item \textbf{Secrets}: éviter tout secret versionné (fichiers \texttt{.env}, mots de passe compose).
  \item \textbf{Cohérence tenant}: s'assurer que toutes les requêtes filtrent bien par \texttt{organization\_id}.
  \item \textbf{Transactions réactives}: clarifier les transaction managers (R2DBC vs JDBC/Liquibase).
  \item \textbf{Observabilité}: compléter Actuator/metrics et corrélation requestId/organizationId.
\end{itemize}

\appendix
\chapter{Annexes}
\chapter{Annexe A --- Spécification API (exhaustive)}
\label{annexe:api}

\section{Introduction}
Cette annexe recense les endpoints \textbf{du module Comptabilité} par contrôleur, tels qu'observés dans le code (\texttt{src/main/java/com/yowyob/erp/accounting/controller}).

\section{\texttt{EcritureComptableController} --- \texttt{/api/accounting/ecritures}}
\begin{itemize}
  \item \textbf{POST} \texttt{/api/accounting/ecritures} --- créer une écriture (\texttt{EcritureComptableDto})
  \item \textbf{PUT} \texttt{/api/accounting/ecritures/\{id\}} --- modifier une écriture (\texttt{EcritureComptableDto})
  \item \textbf{POST} \texttt{/api/accounting/ecritures/\{id\}/validate} --- valider une écriture
  \item \textbf{GET} \texttt{/api/accounting/ecritures} --- lister
  \item \textbf{GET} \texttt{/api/accounting/ecritures/\{id\}} --- obtenir par id
  \item \textbf{GET} \texttt{/api/accounting/ecritures/non-validated} --- lister non validées
  \item \textbf{GET} \texttt{/api/accounting/ecritures/search} --- recherche (start/end/journalId)
  \item \textbf{POST} \texttt{/api/accounting/ecritures/generate} --- générer depuis \texttt{ComptableObject}
  \item \textbf{DELETE} \texttt{/api/accounting/ecritures/\{id\}} --- supprimer
  \item \textbf{PATCH} \texttt{/api/accounting/ecritures/\{id\}/cancel} --- annuler
  \item \textbf{PATCH} \texttt{/api/accounting/ecritures/\{id\}/deactivate} --- désactiver (soft delete)
\end{itemize}

\section{\texttt{CompteController} --- \texttt{/api/accounting/comptes}}
\begin{itemize}
  \item \textbf{POST} \texttt{/api/accounting/comptes} --- créer un compte (\texttt{CompteDto})
  \item \textbf{POST} \texttt{/api/accounting/comptes/generate} --- générer un compte (tiers etc.)
  \item \textbf{GET} \texttt{/api/accounting/comptes} --- lister
  \item \textbf{GET} \texttt{/api/accounting/comptes/\{id\}} --- obtenir par id
  \item \textbf{GET} \texttt{/api/accounting/comptes/search?no\_compte=...} --- recherche
  \item \textbf{GET} \texttt{/api/accounting/comptes/clients} --- comptes clients
  \item \textbf{GET} \texttt{/api/accounting/comptes/suppliers} --- comptes fournisseurs
  \item \textbf{GET} \texttt{/api/accounting/comptes/banks} --- comptes bancaires
  \item \textbf{GET} \texttt{/api/accounting/comptes/cash} --- comptes caisse
  \item \textbf{GET} \texttt{/api/accounting/comptes/type/\{type\}} --- comptes par type
  \item \textbf{PUT} \texttt{/api/accounting/comptes/\{id\}} --- mise à jour
  \item \textbf{DELETE} \texttt{/api/accounting/comptes/\{id\}} --- suppression
\end{itemize}

\section{\texttt{PlanComptableController} --- \texttt{/api/accounting/plan-comptable}}
\begin{itemize}
  \item \textbf{POST} \texttt{/api/accounting/plan-comptable/admin/organizations/\{organizationId\}/plan-comptable/init-ohada-2025} --- initialisation OHADA 2025
  \item \textbf{POST} \texttt{/api/accounting/plan-comptable} --- créer un compte de plan (\texttt{PlanComptableDto})
  \item \textbf{GET} \texttt{/api/accounting/plan-comptable/\{id\}} --- obtenir par id
  \item \textbf{GET} \texttt{/api/accounting/plan-comptable} --- lister
  \item \textbf{GET} \texttt{/api/accounting/plan-comptable/actifs} --- lister actifs
  \item \textbf{GET} \texttt{/api/accounting/plan-comptable/prefix/\{prefix\}} --- filtrer par préfixe
  \item \textbf{GET} \texttt{/api/accounting/plan-comptable/classe/\{classe\}} --- filtrer par classe OHADA
  \item \textbf{PUT} \texttt{/api/accounting/plan-comptable/\{id\}} --- mise à jour
  \item \textbf{POST} \texttt{/api/accounting/plan-comptable/import} --- import (CSV/XLSX multipart)
  \item \textbf{DELETE} \texttt{/api/accounting/plan-comptable/\{id\}} --- désactivation (soft delete)
\end{itemize}

\section{\texttt{JournalComptableController} --- \texttt{/api/accounting/journals}}
\begin{itemize}
  \item \textbf{POST} \texttt{/api/accounting/journals} --- créer (\texttt{JournalComptableDto})
  \item \textbf{PUT} \texttt{/api/accounting/journals/\{id\}} --- mise à jour
  \item \textbf{GET} \texttt{/api/accounting/journals/\{id\}} --- obtenir
  \item \textbf{GET} \texttt{/api/accounting/journals/\{id\}/comptes} --- comptes distincts utilisés
  \item \textbf{GET} \texttt{/api/accounting/journals} --- lister
  \item \textbf{GET} \texttt{/api/accounting/journals/active} --- lister actifs
  \item \textbf{DELETE} \texttt{/api/accounting/journals/\{id\}} --- supprimer
\end{itemize}

\section{\texttt{PeriodeComptableController} --- \texttt{/api/accounting/periodes}}
\begin{itemize}
  \item \textbf{POST} \texttt{/api/accounting/periodes} --- créer (\texttt{PeriodeComptableDto})
  \item \textbf{GET} \texttt{/api/accounting/periodes/\{id\}} --- obtenir
  \item \textbf{GET} \texttt{/api/accounting/periodes} --- lister
  \item \textbf{GET} \texttt{/api/accounting/periodes/code/\{code\}} --- obtenir par code
  \item \textbf{GET} \texttt{/api/accounting/periodes/by-date?date=...} --- obtenir par date
  \item \textbf{GET} \texttt{/api/accounting/periodes/non-closed} --- lister ouvertes
  \item \textbf{GET} \texttt{/api/accounting/periodes/range?start\_date=...\&end\_date=...} --- intervalle
  \item \textbf{PUT} \texttt{/api/accounting/periodes/\{id\}} --- mise à jour
  \item \textbf{PUT} \texttt{/api/accounting/periodes/\{id\}/close} --- clôture
  \item \textbf{DELETE} \texttt{/api/accounting/periodes/\{id\}} --- suppression
\end{itemize}

\section{\texttt{ExerciceComptableController} --- \texttt{/api/accounting/exercices}}
\begin{itemize}
  \item \textbf{POST} \texttt{/api/accounting/exercices} --- créer (\texttt{ExerciceComptableDto})
  \item \textbf{GET} \texttt{/api/accounting/exercices/\{id\}} --- obtenir
  \item \textbf{GET} \texttt{/api/accounting/exercices/\{id\}/periodes} --- périodes de l'exercice
  \item \textbf{GET} \texttt{/api/accounting/exercices} --- lister
  \item \textbf{GET} \texttt{/api/accounting/exercices/active?date=...} --- actif pour date
  \item \textbf{PUT} \texttt{/api/accounting/exercices/\{id\}} --- mise à jour
  \item \textbf{POST} \texttt{/api/accounting/exercices/\{id\}/close} --- clôture exercice
  \item \textbf{DELETE} \texttt{/api/accounting/exercices/\{id\}} --- suppression
  \item \textbf{PATCH} \texttt{/api/accounting/exercices/\{id\}/deactivate} --- désactivation
\end{itemize}

\section{\texttt{RapportController} --- \texttt{/api/accounting/rapport}}
\begin{itemize}
  \item \textbf{GET} \texttt{/api/accounting/rapport/bilan} --- bilan (JSON)
  \item \textbf{GET} \texttt{/api/accounting/rapport/compte-resultat} --- compte de résultat (JSON)
  \item \textbf{GET} \texttt{/api/accounting/rapport/flux-tresorerie} --- flux de trésorerie (JSON)
  \item \textbf{GET} \texttt{/api/accounting/rapport/resume-executif} --- résumé exécutif (JSON)
  \item \textbf{GET} \texttt{/api/accounting/rapport/grand-livre} --- grand livre (JSON)
  \item \textbf{GET} \texttt{/api/accounting/rapport/balance} --- balance (JSON)
  \item \textbf{GET} \texttt{/api/accounting/rapport/balance-agee} --- balance âgée (JSON)
  \item \textbf{GET} \texttt{/api/accounting/rapport/bilan/export/pdf} --- bilan (PDF)
  \item \textbf{GET} \texttt{/api/accounting/rapport/compte-resultat/export/pdf} --- compte résultat (PDF)
  \item \textbf{GET} \texttt{/api/accounting/rapport/grand-livre/export/pdf} --- grand livre (PDF)
  \item \textbf{GET} \texttt{/api/accounting/rapport/balance/export/pdf} --- balance (PDF)
  \item \textbf{GET} \texttt{/api/accounting/rapport/balance-agee/export/pdf} --- balance âgée (PDF)
\end{itemize}

\section{\texttt{PointageController} --- \texttt{/api/accounting/pointage}}
\begin{itemize}
  \item \textbf{POST} \texttt{/api/accounting/pointage/import} --- import relevé bancaire (CSV) + pointage
\end{itemize}

\section{\texttt{LettrageController} --- \texttt{/api/comptable/lettrage}}
\begin{itemize}
  \item \textbf{POST} \texttt{/api/comptable/lettrage/auto} --- lettrage automatique
  \item \textbf{GET} \texttt{/api/comptable/lettrage/status} --- statut (placeholder)
\end{itemize}

\section{Note}
Cette annexe sera étendue au même format pour les autres contrôleurs du domaine (\texttt{taxes}, \texttt{devises}, \texttt{immobilisations}, \texttt{budgets}, \texttt{paie}, \texttt{brouillards}, \texttt{regularisations}, \texttt{rapprochement}, \texttt{releves}, \texttt{notifications}, \texttt{attachments}, etc.).


\chapter{Annexe B --- Dictionnaire de données (exhaustif)}
\label{annexe:data}

\section{Principe}
Le schéma est géré par \textbf{Liquibase}. Une migration structurante (\texttt{changeset-017-unify-tenant-org.xml}) renomme \texttt{tenant\_id} en \texttt{organization\_id} sur la majorité des tables et unifie \texttt{tenants} et \texttt{organizations}.

\section{Tables cœur (comptabilité)}
\subsection{\texttt{organizations}}
\begin{itemize}
  \item \textbf{Rôle}: organisation (tenant) --- unité d'isolation des données.
  \item \textbf{Identifiant}: \texttt{id} (UUID, PK)
  \item \textbf{Attributs} (extrait): \texttt{name}, \texttt{email}, \texttt{telephone}, \texttt{address}, \texttt{tax\_id}, timestamps.
\end{itemize}

\subsection{\texttt{comptes}}
\begin{itemize}
  \item \textbf{PK}: \texttt{id} (UUID)
  \item \textbf{FK}: \texttt{organization\_id} \(\rightarrow\) \texttt{organizations(id)} (après unification)
  \item \textbf{Champs}: \texttt{no\_compte}, \texttt{libelle}, \texttt{classe}, \texttt{type\_compte}, \texttt{solde}, \texttt{actif}, audit.
  \item \textbf{Unicité}: \texttt{(organization\_id, no\_compte)} (contrainte unique).
\end{itemize}

\subsection{\texttt{plans\_comptables}}
\begin{itemize}
  \item \textbf{PK}: \texttt{id} (UUID)
  \item \textbf{FK}: \texttt{organization\_id} \(\rightarrow\) \texttt{organizations(id)}
  \item \textbf{Champs}: \texttt{classe}, \texttt{no\_compte}, \texttt{libelle}, \texttt{notes}, \texttt{actif}, audit.
\end{itemize}

\subsection{\texttt{journaux\_comptables}}
\begin{itemize}
  \item \textbf{PK}: \texttt{id} (UUID)
  \item \textbf{FK}: \texttt{organization\_id} \(\rightarrow\) \texttt{organizations(id)}
  \item \textbf{Champs}: \texttt{code\_journal}, \texttt{libelle}, \texttt{type\_journal}, \texttt{actif}, audit.
\end{itemize}

\subsection{\texttt{exercices\_comptables}}
\begin{itemize}
  \item \textbf{PK}: \texttt{id} (UUID)
  \item \textbf{FK}: \texttt{organization\_id} \(\rightarrow\) \texttt{organizations(id)}
  \item \textbf{Champs}: \texttt{code}, \texttt{libelle}, \texttt{date\_debut}, \texttt{date\_fin}, \texttt{cloture}, audit.
\end{itemize}

\subsection{\texttt{periodes\_comptables}}
\begin{itemize}
  \item \textbf{PK}: \texttt{id} (UUID)
  \item \textbf{FK}: \texttt{organization\_id} \(\rightarrow\) \texttt{organizations(id)}
  \item \textbf{FK}: \texttt{exercice\_id} \(\rightarrow\) \texttt{exercices\_comptables(id)}
  \item \textbf{Champs}: \texttt{code}, \texttt{date\_debut}, \texttt{date\_fin}, \texttt{cloturee}, \texttt{date\_cloture}, audit.
\end{itemize}

\subsection{\texttt{ecritures\_comptables}}
\begin{itemize}
  \item \textbf{PK}: \texttt{id} (UUID)
  \item \textbf{FK}: \texttt{organization\_id} \(\rightarrow\) \texttt{organizations(id)}
  \item \textbf{FK}: \texttt{journal\_id} \(\rightarrow\) \texttt{journaux\_comptables(id)}
  \item \textbf{FK}: \texttt{periode\_id} \(\rightarrow\) \texttt{periodes\_comptables(id)}
  \item \textbf{Champs}: \texttt{numero\_ecriture} (unique), \texttt{libelle}, \texttt{date\_ecriture}, totaux débit/crédit, \texttt{validee}, \texttt{statut}, \texttt{actif}, audit.
\end{itemize}

\subsection{\texttt{details\_ecritures}}
\begin{itemize}
  \item \textbf{PK}: \texttt{id} (UUID)
  \item \textbf{FK}: \texttt{organization\_id} \(\rightarrow\) \texttt{organizations(id)}
  \item \textbf{FK}: \texttt{ecriture\_id} \(\rightarrow\) \texttt{ecritures\_comptables(id)}
  \item \textbf{FK}: \texttt{compte\_id} \(\rightarrow\) \texttt{plans\_comptables(id)} (selon changeset initial)
  \item \textbf{Champs}: \texttt{sens}, \texttt{montant\_debit}, \texttt{montant\_credit}, \texttt{lettree}, \texttt{date\_lettrage}, \texttt{pointee}, \texttt{reference\_bancaire}, audit.
\end{itemize}

\subsection{\texttt{journal\_audit}}
\begin{itemize}
  \item \textbf{PK}: \texttt{journal\_audit\_id} (UUID) / entité \texttt{JournalAudit}
  \item \textbf{FK}: \texttt{organization\_id} \(\rightarrow\) \texttt{organizations(id)}
  \item \textbf{FK}: \texttt{ecriture\_id} \(\rightarrow\) \texttt{ecritures\_comptables(id)} (SET NULL)
  \item \textbf{Champs}: \texttt{action}, \texttt{date\_action}, \texttt{utilisateur}, \texttt{details}, avant/après, audit.
\end{itemize}

\section{Tables avancées (extraits)}
\subsection{\texttt{brouillards\_comptables}}
Table de brouillard/draft (source externe, pièce, JSONB \texttt{data\_json}, statut, lien optionnel vers \texttt{ecriture\_id}).

\subsection{\texttt{immobilisations} / \texttt{amortissements\_lignes}}
Gestion des immobilisations et amortissements, avec liens vers des comptes (\texttt{compte\_immo\_id}, \texttt{compte\_amort\_id}, \texttt{compte\_dotation\_id}) et exercice.

\section{Note de complétude}
Ce dictionnaire est aligné sur les changesets observés (\texttt{001..020}). Il peut être enrichi section par section avec:
\begin{itemize}
  \item types exacts, nullabilité et valeurs par défaut,
  \item index/contraintes (unique, FK, cascade),
  \item règles métier associées par service.
\end{itemize}


\section{Principaux fichiers de référence}
\begin{itemize}
  \item \texttt{pom.xml}
  \item \texttt{src/main/resources/application.properties}
  \item \texttt{src/main/resources/db/changelog/db.changelog-master.xml}
  \item \texttt{src/main/java/com/yowyob/erp/config/auth/SecurityConfig.java}
  \item \texttt{src/main/java/com/yowyob/erp/config/organization/OrganizationWebFilter.java}
\end{itemize}

\end{document}

